\documentclass[10pt,a4paper]{article}

\usepackage[T1]{fontenc}
\usepackage[utf8]{inputenc}
\usepackage{lmodern}
\usepackage[margin=1.18cm,top=1.12cm,bottom=1.25cm]{geometry}
\usepackage{microtype}
\usepackage{enumitem}
\usepackage{array}
\usepackage{fancyhdr}
\usepackage[hidelinks]{hyperref}

\setlength{\parindent}{0pt}
\setlength{\parskip}{1.8pt}
\setlist[itemize]{leftmargin=*,nosep}
\renewcommand{\arraystretch}{0.94}
\setlength{\tabcolsep}{3pt}
\setcounter{secnumdepth}{0}

\pagestyle{fancy}
\fancyhf{}
\cfoot{\footnotesize \thepage\ of 2}
\renewcommand{\headrulewidth}{0pt}

\newcommand{\member}[2]{%
  \vspace{2pt}%
  {\bfseries\large #1}\hfill{\small #2}\par
  \vspace{1pt}\hrule\vspace{3pt}%
}
\newcommand{\role}{\textbf{Member responsibility.} }
\newenvironment{timeline}{%
  \vspace{2pt}%
  \begin{tabular}{@{}>{\bfseries}p{0.12\linewidth}p{0.83\linewidth}@{}}%
  \hline
  Week & \textbf{Planned work}\\
  \hline
}{%
  \hline
  \end{tabular}%
  \vspace{2pt}%
}
\newcommand{\wk}[2]{#1 & #2\\}

\begin{document}

\begin{center}
  {\LARGE\bfseries Interactive FPGA Graphical Calculator}\par
  \vspace{2pt}
  {\large GPT-6 Project Proposal --- Group ID 06}
\end{center}

\vspace{-2pt}
\begin{minipage}{\textwidth}
    \textbf{Group theme.}
    This project will implement an interactive graphical calculator as one
    integrated Verilog design on a Basys3 FPGA with a $96\times64$ OLED display.
    A unified main menu will let users enter signed values with the board
    switches and select either a conventional calculator or Graph Mode. The
    calculator will perform addition, subtraction, multiplication and division,
    while detecting invalid operations and overflow and retaining recent
    results. Graph Mode will accept one or two functions through a guided
    coefficient-entry sequence, render coloured linear or quadratic graphs on
    scaled axes, and calculate the intersection of two linear functions. Clear
    status messages will distinguish parallel, identical, invalid and off-screen
    cases. Modular controller, arithmetic, input and rendering interfaces will
    support independent testing before final integration. Extension candidates
    include automatic view fitting,
    zoom, pan, a movable coordinate cursor, graph tracing, saved function
    presets, cubic, absolute-value, reciprocal, square-root, step, piecewise and
    LUT-based sine/cosine graph rendering, a coordinate-value table, automatic
    root, intercept and turning-point markers, numerical slope readout, memory
    recall and selectable grid density.
\end{minipage}

\vspace{5pt}

\member{Milla Gl\"annfjord}{Group Member 1}
\role
The main responsibility of this group member is to develop a module that reads
switches and convert their states into binary numbers, including support for
signed values. This group member will also write one module for detecting
invalid user inputs and one potential overflow, and implementing a history-data
module that stores and supplies the three most recent results to Hanna's
calculator interface. The modules will be designed so that the other group
members can use it for both the calculator and graphing functions.

\begin{timeline}
\wk{7}{Complete the practice task provided in the project manual to gain a
better understanding of the OLED, Verilog programming and the different
components needed for the project.}

\wk{8}{Develop the module responsible for reading the switches and interpreting
their states as binary numbers. Ensure that the module can handle both positive
and negative values and that it can be integrated with the other group members'
modules.}

\wk{9}{Implement additional functionality to detect invalid user inputs and
potential overflow during calculations. Develop a history module that stores
the three most recent calculation results.}

\wk{10}{Have a group meeting to review the progress of the project and to decide
how to divide the remaining graphical components between the group members.
Discuss how the different modules will be connected and integrated into the
complete system.}

\wk{11}{Develop the signed-input preview and history/error panel for the OLED
display. Test these graphical components and adjust their interfaces so that
they work correctly with the calculator module.}

\wk{12}{Have a final group meeting to integrate all the modules and complete the
project. Perform final testing, identify and fix any remaining issues, and
ensure that the system meets the project requirements.}
\end{timeline}

\member{Hanna Wahlgren}{Group Member 2}
\role
The main responsibility is to implement the calculator module,
supporting the four basic arithmetic operations: addition, subtraction,
multiplication and division. The calculator will use the FPGA board's buttons
to select operations and switches to enter signed binary numbers. Another
responsibility is to develop the calculator interface on the OLED display,
including digits, mathematical symbols, input values and calculation results.
It will also develop the main menu logic in
the top module, allowing users to navigate between the calculator and graphing
functions, select a program and return to the main menu.

\begin{timeline}
\wk{7}{Completed the practice tasks in the project manual to gain a better
understanding of the OLED display and overall Verilog programming required for
the project.}

\wk{8}{Started developing the calculator module and implemented the basic
arithmetic operations using a state-based approach. Created a function for
drawing digits (0--9) on the OLED display and began implementing the graphical
representation of input values, mathematical operators and calculation results.}

\wk{9}{Developed the main-menu logic in the top module, allowing users to
navigate between the calculator and graphing programs using the buttons.
Implemented functionality to select a program and return to the main menu.
Worked on integrating Milla's input, validation, overflow and history-data
modules into the calculator interface.}

\wk{10}{Test the calculator module and its integration with the other modules.
Improve the OLED graphics, including negative numbers, results and error
messages. Participate in the group meeting to divide the remaining graphical
tasks.}

\wk{11}{Implement the assigned graphical components, including the main-menu
interface. Improve the calculator's graphical layout and test navigation
between the menu and the different programs.}

\wk{12}{Integrate the calculator and menu with the other group members' modules.
Perform final testing on the FPGA board, fix remaining issues and prepare the
system for the final demonstration.}
\end{timeline}

\newpage

\member{Wang Qichen}{Group Member 3}
\role
Design and implement the graph-rendering component of our FPGA-based graphical
calculator. The renderer will convert the OLED \texttt{pixel\_index} into
coordinates for the $96\times64$ display and draw a readable graph area with
axes, grid lines and labels. It will display two signed linear functions,
$y=k_1x+m_1$ and $y=k_2x+m_2$, using different colours and continuous
line-drawing logic. Using graph-space results supplied by the controller, it
will draw intersection and feature markers and display status messages for
parallel, identical, invalid or off-screen cases. During coefficient entry, the
OLED will show the input stage and instructions for the board
controls.

\begin{timeline}
\wk{7}{Completed the OLED familiarisation tasks and studied RGB565 colours,
\texttt{pixel\_index} conversion and basic Verilog drawing techniques.}

\wk{8}{Investigated coordinate scaling, signed graph values, line-rendering
methods and the interface required between the graph controller and OLED
renderer.}

\wk{9}{Defined the renderer architecture, colour scheme, graph states,
intersection display and exceptional-case overlays, and coordinated the
required signals with the controller component.}

\wk{10}{Implement and simulate the core axes, grid, two-line renderer,
input-stage text and intersection marker. Add an equation legend and
graph-visibility controls, then integrate the OLED renderer.}

\wk{11}{Render zoom, pan and cursor states supplied by the controller, with graph
clipping, automatic scale selection, coordinate labels and cursor snapping.
Implement a coordinate-value table and selectable grid density, and display
controller-supplied root, intercept, turning-point and slope results. Prototype
prioritised rendering for quadratic, cubic,
absolute-value, reciprocal, square-root, step, piecewise and LUT-based
sine/cosine functions. Add
adaptive tick labels, asymptote and discontinuity markers, a split graph/table
view and parameter animation.}

\wk{12}{Optimise resource usage, verify timing, test signed and extreme inputs
on the Basys3, complete group integration and prepare a reliable demonstration.
Add saved display presets, a guided demo screen and clearer
on-screen help for all graph controls.}
\end{timeline}

\member{Oskar Plogner}{Group Member 4}
\role
The main responsibility is to develop the graph controller and mathematical
calculations for the graphing mode. This includes managing user inputs through
a finite state machine (FSM) for selecting one or two graphs, function types and
signed coefficients. Another major responsibility is
calculating intersections between linear functions using fixed-point
arithmetic, including handling parallel or identical lines, division by zero
and overflow. The controller will communicate with the graph renderer by
providing coefficients, calculation results and control signals. The controller
will also own zoom, pan and cursor state.

\begin{timeline}
\wk{7}{Completed the OLED familiarisation tasks and discussed possible project
functionalities and responsibilities within the group.}

\wk{8}{Planned the graphing functionality, including coefficient input, graph
selection, intersection calculations and communication between the controller
and renderer.}

\wk{9}{Developed the initial \texttt{graph\_controller.v} module using an FSM.
Implemented selection of one or two graphs, linear/quadratic function selection,
coefficient storage, reset functionality and signals indicating when the graphs
are ready to be plotted. The module has not yet been tested or integrated.}

\wk{10}{Simulate and test the graph controller with different input
combinations. Begin integration with the signed-number input module and OLED
renderer, adapting the interfaces between the modules. Implement intersection
calculations for two linear functions, including handling parallel and
identical lines.}

\wk{11}{Implement error handling for invalid calculations, overflow and
intersections outside the visible graph area. Continue testing and improving
integration with the other modules. Implement and output zoom, pan and cursor
state for Wang's renderer. Add function selection, parameters and fixed-point
evaluation for prioritised non-linear graphs, together with quadratic/cubic
intersection solving, tangent/normal calculations, area estimation,
discontinuity detection and parameter-animation control.}

\wk{12}{Integrate the graph controller and intersection calculations with the
other modules. Test the complete graphing functionality on the Basys3 board,
fix remaining issues and help prepare the final demonstration.}
\end{timeline}

\end{document}
